\chapter*{A reviewer's route}
\addcontentsline{toc}{chapter}{A reviewer's route}

These pages are for a specialist who has been asked to referee the paper and wants the shortest path through it: the
statement, the chain of the proof with the place where each link is proved, the joints to check in order, what is
imported, and what the numbers rest on. \Cref{ch:paper} reads the paper section by section; \Cref{ch:congruence} holds the
proofs of Section 4. Paper labels are quoted as text (Proposition prop:integral), as in the rest of the book.

\section*{1. The statement}

\emph{Hypotheses.} $E/\Q$ is the $X_0(N)$-optimal curve of its isogeny class, $f=f_E$; $\ell$ is an odd prime with
$\ell\nmid2N$, and $\rhobar_{E,\ell}$ is absolutely irreducible on $\Gal(\Qbar/\Q(\sqrt{\ell^*}))$, $\ell^*=(-1)^{(\ell-1)/2}\ell$;
that is, $\ell\notin S_f$ in the notation of Diamond, Flach and Guo. For $\ell\nmid2N$ this says: $E[\ell]$ is irreducible
and, if $\ell=3$, the mod $3$ image is not $N_s(\F_3)$ (Proposition prop:Sf). With $T=\Tl E$, Theorem thm:main states
\begin{gather*}
v_\ell(\degphi)=\len_{\Zl}\Hf(\Q,\ad T\otimes\Qlmod)+\sum_{q\mid N}v_\ell\bigl(\Tam_q(\ad T)\bigr)+\sum_{q\mid N}e_\ell(E,q),
\tag{eq:identity}\\
v_\ell(\degphi)\ \ge\ \sum_{q\mid N}e_\ell(E,q)+\sum_{q\mid N}t_\ell(E,q),\tag{eq:thm}
\end{gather*}
where $e_\ell(E,q)=v_\ell\bigl(P_q(q^{-2})/(1-a_q(E)^2q^{-2})\bigr)$ is the valuation of the Euler factor of $L(\ad f,s)$ at
$s=1$ that the naive adjoint $L$-function omits (zero unless $q$ is additive), and $t_\ell(E,q)=v_\ell(-v_q(j_E))$ when
$v_q(j_E)<0$, and $0$ otherwise (Definition def:terms). The second line follows from the first because the Selmer length is
nonnegative and $v_\ell(\Tam_q(\ad T))\ge t_\ell(E,q)$ (Lemma lem:tamfactor). Conjecture conj:euler, which extends the
inequality to $\ell\mid N$ and to the image $\Ns$, is not part of the theorem.

\section*{2. The chain of the proof}

The paper's proof of Theorem thm:main is the table at the start of its Section 4.3; the rows below follow it. ``Paper''
means proved in the paper, ``book'' that the book gives a complete proof (expanding the paper's), ``imported'' that the
statement is cited.

\begin{center}\footnotesize
\begin{tabular}{@{}>{\raggedright\arraybackslash}p{0.33\textwidth}>{\raggedright\arraybackslash}p{0.17\textwidth}%
>{\raggedright\arraybackslash}p{0.18\textwidth}>{\raggedright\arraybackslash}p{0.24\textwidth}@{}}
\toprule
step & paper label & in the book & status\\
\midrule
$H^1(X_0(N),\Zl)$ is the integral structure of DFG, Hecke equivariantly, $U_p$ included & prop:integral &
\Cref{ch:congruence} & paper; book (full proof)\\
$v_\lambda(\eta^\emptyset_{f,\lambda})=v_\ell(\degphi)$ (projection formula) & prop:pullback, lem:hecke &
\Cref{ch:moddeg,ch:congruence} & paper; book; uses \cite[Theorem 2.2]{ARS12} for $r_E$\\
$\eta^\Sigma=\eta^\emptyset\prod_{p\in\Sigma}L_p^{\nv}(A_f,1)^{-1}$ & thm:dfg14 & \Cref{ch:congruence} &
imported, \cite[Proposition 1.4(c)]{DFG04}\\
$\ell\notin S_f$ read for elliptic curves; $H^0(\Q,W^D)$ of length $h_\ell(E)$ & prop:Sf, lem:h0, prop:ns &
\Cref{ch:galrep,ch:reps} & paper; book\\
$v_\lambda(\eta^\Sigma)=\len\HSigma(\Q,W)$ & thm:dfg37 & \Cref{ch:congruence} (outline) &
imported, \cite[Theorem 3.7]{DFG04}\\
$H^0(\Q,V^D)=\Hf(\Q,V^D)=0$ & thm:dfg27 & \Cref{ch:blochkato,ch:congruence} & imported, \cite[Theorem 2.7]{DFG04}\\
$\HSigma(\Q,W)$ against $\Hf(\Q,W)$ & lem:dfg21 & \Cref{ch:blochkato,ch:galcoh} &
imported, \cite[Lemma 2.1]{DFG04}; the paragraph after it paper, book\\
$\Fitt\Hf(\Q_p,T^D)=L_p(A_f,1)^{-1}\Tam_p(A)$ & lem:fpr & \Cref{ch:blochkato,ch:lfunctions} &
imported, \cite[I.4.2.2]{FPR94}, \cite[(57)]{DFG04}; first identity also proved in book\\
values of $e_\ell$; shapes of $P_q$ & def:terms, lem:euler, prop:shapes & \Cref{ch:lfunctions,ch:galrep,ch:reduction} &
paper; book\\
$\Tam_q(\ad T)$ at the bad primes & lem:tamfactor, lem:tamadd & \Cref{ch:blochkato,ch:galcoh} & paper; book\\
the assembly & proof of thm:main, cor:explicit & \Cref{ch:congruence} & paper; book\\
\bottomrule
\end{tabular}
\end{center}

\section*{3. The delicate joints, in order}

\begin{enumerate}[label=(\alph*), leftmargin=*]
\item \emph{Proposition prop:integral.} This is the one input that the paper says is ``not a published statement in the
form we need''; the paper proves it, and adds that ``it holds integrally, with no localisation.'' The proof identifies
DFG's Betti lattice $M(M,1)_{!,B}$, defined as invariants of the cohomology of a modular curve of full level, with
$H^1(X_0(M),\Z)$: compactly supported Hochschild-Serre and the equivariant Leray sequence of the finite quotient map (Step 1),
the image of $H^1_c\to H^1$ (Step 2), Hecke operators checked over $\C$ (Step 3), base change to $\Zl$ (Step 4). To check: the
reading of DFG's Sections 1.2 and 1.4.1 (Lemma 1.1, Proposition 1.3 for the Hecke action with $U_p$, display (5)), and the
paragraph after the proposition, where the pairing $\hat\delta_!$ is $\frac u2$ times the cup product composed with $w_{M}$,
$u$ an $\ell$-adic unit, by the statement before DFG's Lemma 1.2. Only the unit matters for odd $\ell$; the exact constant is
not proved. \Cref{ch:congruence} gives the full proof; DFG's sections are cited there by section, as the paper cites them.
\item \emph{The DFG hypothesis chain.} Theorem thm:dfg37 needs $\ell\notin S_f$ and $\Sigma$ the set of primes dividing $N$
(minimal ramification outside $\Sigma$). Proposition 1.4(c) needs $\ell\notin\Sigma$ and $\rhobar$ irreducible; the paper
shows by \curve{11a} at $\ell=3$ that $\ell\notin\Sigma$ cannot be dropped. Lemma 2.1 needs $H^0(\Q,V^D)=\Hf(\Q,V^D)=0$:
the first is automatic, as $H^0(\Q_p,V^D)=0$ for $p\in\Sigma$; the second is Theorem 2.7; and its local terms are the torsion of
$H^1(\Q_p,T^D)$ because $p\ne\ell$. The last term $H^0(\Q,W^D)^*$ of the sequence is dropped by Lemma lem:h0, which gives
length $h_\ell(E)=0$ exactly when $\ell\notin S_f$. A referee should confirm that the four DFG statements are used in the
published form and numbering (\Cref{ch:paper} lists which were compared with the published text).
\item \emph{The local terms at $q=2,3$.} In (eq:identity) the terms at $2$ and $3$ are those of Definition def:terms with
PARI's Euler factor; the paper states that \texttt{lfunsympow(E,2)} ``returns these factors at every prime, including $2$
and $3$''. The identity needs only the definition; Proposition prop:shapes (four shapes of $P_q$ at every bad prime, proved
from the image of inertia) makes them explicit, and Remark rem:wild gives the examples \curve{336f1}, \curve{1232l1},
\curve{592c1}. The file \texttt{final/pipeline/samples/wild\_euler\_q23.csv} (with its \texttt{README.md} and log) records
$P_q$, $P_q(q^{-2})$ and $e_\ell$ for $3{,}305$ rows at $q=2,3$ ($2{,}905$ curves: those named in the paper and the optimal
non-CM curves with $N<2{,}000$ and $4\mid N$ or $9\mid N$), each matched to one of the four shapes, with no mismatch.
\item \emph{The image $\Ns$ and $h_3$.} Theorem thm:main excludes these curves ($3\in S_f$). The paper's ``Status of the
correction'' separates three things. Proved: $h_\ell(E)\le1$, with equality exactly when $\ell=3$ and the mod $3$ image is
$N_s(\F_3)$ (Lemma lem:h0, Proposition prop:ns). Conditional: that the correction is the term $-h_\ell(E)$ that Lemma 2.1
produces assumes $H^1_f(\Q,V^D)=0$, ``which is not known for these curves'', and that $v_\lambda(\eta^\Sigma_{f,\lambda})$ is
still the Selmer length (Remark rem:ns). Empirical: the weights $w_3$ at $q=\ell=3$ and the shortfall pattern, a failure of
(eq:thm) by exactly one on $364$, $68$ and $43$ curves of the three table ranges, all with image $\Ns$
(\Cref{ch:galrep,ch:blochkato}).
\item \emph{Theorem thm:quant, Step 2.} The proof is a chain of equalities with one inequality: Lemma lem:ddt417 (the
inequality half of \cite[Lemma 4.17]{DDT}, proved in the paper) applied to $H^1(X_0(N),\Zl)_{\mf m}$ and the depleted form
$g^S$ of the twist, giving $v_\ell(\eta^{\mathrm{Hecke}}(g^S))\ge v_\lambda(\eta^S_{g,\lambda})$; in the paper's words, ``the
only inequality is Step 2''. To check: that $L[\wp]$ is the lattice $M_{g^S,\lambda}$ (Proposition prop:integral at level
$N$), that the pairing is perfect on the $\mf m$-component ($e_{\mf m}$ self-adjoint), and that freeness is not used. The
inequality needs only $E[\ell]$ irreducible and $\ell\nmid2N$; the equality uses Theorem thm:main for $E$ and $E'_S$
(\Cref{ch:congruence}).
\end{enumerate}

\section*{4. Imported results}

As the paper cites them, for Theorem thm:main and Theorem thm:quant:
\cite[Proposition 1.4(c)]{DFG04};
\cite[Theorem 3.7]{DFG04} (the form given as Theorem 0.2 there);
\cite[Theorem 2.7]{DFG04};
\cite[Lemma 2.1]{DFG04}, with $S_f$ from \cite[(31)]{DFG04};
\cite[I.4.2.2]{FPR94}, as recalled in the proof of \cite[Theorem 2.15]{DFG04}, with $\Tam_p(T^D)=\Tam_p(A)$ from \cite[(57)]{DFG04};
for Proposition prop:integral and the pairing, \cite[\S\S1.2, 1.4]{DFG04}, \cite[Lemma 1.1]{DFG04}, \cite[\S\S1.2.1, 1.3.2]{DFG04} and
\cite[\S1.5.2, before Lemma 1.2]{DFG04};
\cite[Theorem 2.2]{ARS12} ($\degphi\mid r_E$, with equal $p$-parts for $p^2\nmid N$);
\cite[Lemma 4.17]{DDT} (its inequality half is reproved);
\cite[\S1.11, Propositions 11 and 12]{Serre72} (inertia at $\ell$, in Proposition prop:Sf);
\cite[\S2, Theorem 2]{SerreTate68} and \cite{Kraus90} (inertia images at additive primes);
Tate's algorithm \cite{Tate75}, \cite[IV.9]{SilvermanATAEC};
and PARI's \texttt{lfunsympow} \cite{PARI} for the factors at $2$ and $3$. The status of each citation (compared with the
published text, or as cited) is the last table of \Cref{ch:paper}.

\section*{5. The computations}

No computation enters the proof of Theorem thm:main or Theorem thm:quant. The data concern Conjectures conj:euler and
conj:eis, the sharpness of the hypothesis, and checks of the local lemmas. They rest on three kinds of evidence.
\begin{itemize}[leftmargin=*]
\item \emph{Large ranges} (conductor up to $500000$, and the $893$ curves outside the tables): computed by the scripts in
\texttt{final/pipeline}. The pipeline was rerun in full in this repository from the pinned data, with PARI/GP 2.17.2
(\texttt{scripts/\allowbreak repro\_run.sh}, then \texttt{scripts/\allowbreak reproduce\_tables.sh}), and reproduced every committed result
summary and every table file byte for byte, the four twisted rows of Table tab:isolated included; the out-of-table curves
were not drawn again, their table being regenerated from the committed curve files. A third stage
(\texttt{scripts/stage3\_*.py}) recomputed the numbers that the paper states without a committed result file: $21$ of $23$
agree as stated, and two sentences of the paper were corrected in wording (Section sec:two: $19$ exceptions on the
extension, not $20$; Section sec:eis: the $18$ curves with a point of order $9$ are all such curves, optimal or not). The
logs, the timings and two caveats (\texttt{galrep} labels written without dots above conductor $400000$, which only weaken
one term of the test of Conjecture conj:eis, and a different number of shards) are in the pipeline's \texttt{README.md},
section ``Full rerun''; \Cref{ch:computation} summarises them. The printed tables were also compared with the committed
result files (\texttt{scripts/selftest.py}, output in \texttt{results/selftest.log}; \texttt{MANIFEST.sha256} fixes every
file). The book's own scripts do not reach this range.
\item \emph{Conductor below $10{,}000$} ($38{,}042$ optimal curves): recomputed for the book with PARI/GP 2.17.2 from the
same pinned commit, including Table tab:excess (\Cref{ch:blochkato}), the check of Lemma lem:euler against PARI's factors,
and the counts of Sections sec:tam, sec:anomaly and sec:wild that fall in this range (\Cref{ch:computation}). These runs were
made with the paper in hand: they are checks, not blind tests.
\item \emph{Single curves}: the worked verifications of \Cref{ch:paper} and the function \texttt{rule} of the paper's
appendix (\texttt{scripts/euler\_rule.gp}).
\end{itemize}
Which ranges were used to formulate each statement and which to test it is set out, statement by statement, in the
protocol remarks of \Cref{ch:computation}. Each computational claim of the paper,
with the evidence behind it, is listed in the claim ledger in \texttt{final/claims/}.

\section*{6. What has not been checked}

The paper and this book were produced with an AI system under the author's direction, and neither has yet been read line
by line by a human expert; the paper's statement on the use of an AI system and the acknowledgement in the preface say so.
